\begin{helper}[The abstract comparison scheme: Campanato decay]
Let $n\ge2$, $p>1$, $0<\alpha\le1$, $\delta'>0$, $c_1,c_3>0$ and
\[
0<\tau<p\alpha\theta,\qquad \theta=\frac{\delta'}{n+p\alpha} .
\]
Then (A) and (B) below both hold.

\medskip
\noindent\textbf{(A) The scheme, from a comparison bundle.}  There are
$\lambda_0\in(0,p)$ and $c_{10},R_2>0$ with $R_2\le1$, depending only on
$n,p,\alpha,\delta',\tau,c_1,c_3$, with the following property.  (The
upper bound on $\lambda_0$ is \emph{strict}; this is what the consumers need,
because Campanato's theorem (C5) is available at the exponent
$\lambda=n+\lambda_0$ only for $\lambda<n+p$.  The paper's Theorems 4 and 6
state $0<\lambda_0\le p$, which is weaker and which they keep.)

Let $0\le\varepsilon\le1$, let $\Omega\subseteq\R^n$ be open, let $\Omega'$ be
open with $\overline{\Omega'}\subseteq\Omega$, let $R_1>0$ and let
$u\in W^{1,p}(\Omega)$ with weak gradient $g=\nabla u$ (\noderef{MemW1pOn}).
Write
\[
\Lambda(x_0,R)=R^{-n}\int_{B_R(x_0)}(1+|\nabla u|^2)^{\frac p2}dx .
\]
Assume
\begin{align}
\Lambda(x_0,R)&\le c_1R^{-\tau}\qquad (x_0\in\Omega',\ 0<R\le R_1),\tag{4.15}
\end{align}
and that for every $x_0\in\Omega'$ and $0<R\le R_1$ with
$\overline B_R(x_0)\subseteq\Omega$ there is a comparison map
$v\in W^{1,p}\bigl(B_R(x_0)\bigr)$ with weak gradient $g_v$
(\noderef{MemW1pOn}) such that
\begin{align}
\int_{B_\rho(x_0)}\bigl|g_v-(g_v)_{B_\rho(x_0)}\bigr|^p dx
&\le c_3\Bigl(\frac\rho R\Bigr)^{p\alpha}\rho^{\,n}\Lambda(x_0,R)
\qquad(0<\rho\le R/2),\tag{4.16}\\
\int_{B_R(x_0)}|\nabla u-g_v|^p\,dx&\le R^{\delta'+n}\Lambda(x_0,R).\tag{4.17}
\end{align}
Then for every $x_0\in\Omega'$, every $R$ with $0<R\le\min\{R_1,R_2\}$ and
$\overline B_R(x_0)\subseteq\Omega$, and every $\rho>0$ with
$2\rho\le R^{1+\theta}$,
\[
\int_{B_\rho(x_0)}\bigl|\nabla u-(\nabla u)_{B_\rho(x_0)}\bigr|^p dx
\le c_{10}\,\rho^{\,n+\lambda_0}.
\tag{4.19}
\]
Here $(h)_{B_\rho(x_0)}=|B_\rho|^{-1}\int_{B_\rho(x_0)}h\,dx$ is the vector mean.

\noindent\textbf{On the range of $\rho$.}  The conclusion is stated for the
radii $\rho$ with $2\rho\le R^{1+\theta}$ at an admissible comparison radius $R$,
and not for all $\rho\le\min\{R_1,R_2\}$, because that is exactly the range the
hypotheses reach.  The thesis derives (4.19) of the thesis by putting
$\rho=\frac12R^{1+\theta}$ (p.~45), so the estimate at a given $\rho$ consumes
the comparison bundle at the \emph{larger} radius $R=(2\rho)^{1/(1+\theta)}$ ---
larger because $\theta>0$ and $2\rho\le1$ --- and the bundle is available only
for radii $\le R_1$ whose closed ball lies in $\Omega$.  A conclusion quantified
over all $\rho\le\min\{R_1,R_2\}$ would assert the decay at radii whose
associated comparison radius exceeds $R_1$, where the hypotheses say nothing.
Conversely nothing is lost: as $R'$ runs over $(0,R]$ the number
$\frac12R'^{\,1+\theta}$ runs over all of $(0,\frac12R^{1+\theta}]$, so every
$\rho$ in the stated range is of the thesis' special form for exactly one
$R'\le R$, and the containment $\overline B_R(x_0)\subseteq\Omega$ passes to
$\overline B_{R'}(x_0)$.  A consumer fixes $R$ from the geometry of its
subdomain once and then has (4.19) of the thesis for all small $\rho$, which is what
Campanato's characterisation of the H\"older classes consumes.

\medskip
\noindent\textbf{(B) The scheme, from the producer-side data.}  Let $r>0$ and
$0<\alpha_H\le1$, suppose the Sobolev--Poincar\'e inequality on balls holds in
dimension $n$ (\noderef{CitedSobolevPoincareBall}), and assume the cited
solvability result (C2) (\noderef{CitedDirichletW1p}), the cited regularity
result (C1) (\noderef{CitedRegularityEpsPos}) and the cited Dirichlet result
(C3) with H\"older boundary data (\noderef{CitedDirichletHolder}), the last
three for \emph{every} field $A'\colon\R^n\to\R^n$ and all parameters
$\varepsilon',\gamma_0',\gamma_1'$ with $\gamma_0'>0$, $\gamma_1'>0$ --- and,
for (C1), for every open $\Omega_0\subseteq\R^n$.  Then there is
$\alpha_2\in(0,1]$, depending only on $n,p,r$ and $\alpha_H$, such that for all
$\delta_2>0$ and $\tau_2$ with
\[
0<\tau_2<p\alpha_2\theta_2,\qquad \theta_2=\frac{\delta_2}{n+p\alpha_2},
\]
there are $\lambda_2\in(0,p)$ and $R_3\in(0,1]$, depending only on
$n,p,r,\alpha_H,\delta_2,\tau_2$, with the following property: for every
$c_2>0$ there is $c_{20}>0$ such that the following holds.

Let $0\le\varepsilon\le1$, let $\gamma_0,\gamma_1>0$ with
$\gamma_1\le r\gamma_0$, let $\Omega\subseteq\R^n$ be open, let $\Omega'$ be
open with $\overline{\Omega'}\subseteq\Omega$, let $R_1>0$, let
$\mathcal A\colon\R^n\times\R\to(\R^n\to\R^n)$ be a family of fields, and let
$u,U\in W^{1,p}(\Omega)$ both have weak gradient $g=\nabla u$
(\noderef{MemW1pOn}).  Assume
\begin{itemize}
\item $\overline B_R(x_0)\subseteq\Omega$ for every $x_0\in\Omega'$ and every
  $0<R\le R_1$;
\item for every $x_0\in\Omega'$ the normalised frozen field
  $h\mapsto\mathcal A(x_0,U(x_0))(h)-\mathcal A(x_0,U(x_0))(0)$ satisfies
  (Q1) of the thesis--(Q2) of the thesis on $\R^n$ with the parameters
  $p,\varepsilon,\gamma_0,\gamma_1$ (\noderef{EllipticStructureQ12});
\item $U\in C^{\alpha_H}\bigl(\overline B_{R_1}(x_0)\bigr)$
  (\noderef{MemHolderOn}) for every $x_0\in\Omega'$;
\item the Morrey bound $\Lambda(x_0,R)\le c_2R^{-\tau_2}$ for $x_0\in\Omega'$
  and $0<R\le R_1$;
\item the smallness (4.17) of the thesis, quantified over bundles: for every
  $x_0\in\Omega'$, every $0<R\le R_1$ and every $v,V,g_v$ for which the frozen
  comparison bundle
  $\mathrm{FrozenComparisonOn}(n,p,x_0,R,\mathcal A(x_0,U(x_0)),U,U,g,v,V,g_v)$
  holds (\noderef{FrozenComparisonOn}),
  \[
  \int_{B_R(x_0)}|\nabla u-g_v|^p\,dx\le R^{\delta_2+n}\Lambda(x_0,R).
  \]
\end{itemize}
Then for every $x_0\in\Omega'$, every $R$ with $0<R\le\min\{R_1,R_3\}$ and every
$\rho>0$ with $2\rho\le R^{1+\theta_2}$,
\[
\int_{B_\rho(x_0)}\bigl|\nabla u-(\nabla u)_{B_\rho(x_0)}\bigr|^p dx
\le c_{20}\,\rho^{\,n+\lambda_2}.
\]

\medskip
\noindent\textbf{What (B) adds, and why its exponent is produced rather than
received.}  Clause (A) asks its consumer for a comparison bundle: a $g_v$
satisfying (4.16) of the thesis and (4.17) of the thesis at every centre and radius.
Producing that bundle is the same work in \S4 and in \S5 --- solve the frozen
Dirichlet problem, then read (4.16) of the thesis off Theorem~Theorem~2 of the thesis and
Theorem~Theorem~1 of the thesis --- and clause (B) does it once, here, so that
Theorem~Theorem~4 of the thesis and Theorem~Theorem~6 of the thesis supply only what is genuinely
theirs: the structure of their own frozen field, the H\"older datum for the
representative, the Morrey bound, and (4.17) of the thesis, which is their own
Lemma~Lemma~4.1 of the thesis resp.\ Lemma~Lemma~5.1 of the thesis.  Clause (A) is unchanged and
Theorems 3 and 5 are untouched.

The exponent $\alpha_2$ is \emph{produced} by (B) and not received, because it
is Theorem~Theorem~2 of the thesis's oscillation exponent, which that theorem binds in its
own existential after $n,p$ and the ratio $r$; a consumer cannot name it in
advance.  That is why $\delta_2$, $\tau_2$ and $c_2$ are bound after it, in the
order the thesis' own proof of Theorem~Theorem~4 of the thesis uses: first the exponent,
then ``given $\tau>0$'' and the Morrey constant $c_1$ that Theorem~3 supplies at
that $\tau$, and only then the admissibility $\tau<p\alpha\theta$.  The
H\"older exponent $\alpha_H$ of the representative is a separate, \emph{bound}
parameter, fixed by the consumer at $\alpha_H=\tfrac12$ from part (iii) of
Theorem~Theorem~3 of the thesis resp.\ Theorem~Theorem~5 of the thesis; it must not be taken from (C4),
whose exponent depends on $M=\|u\|_{L^\infty}$ while $\lambda_2$ may not.

\medskip
\noindent\textbf{On the quantifier order of (B).}  $\lambda_2$ is bound before
$c_2$ and before all the data, so it depends only on $n,p,r,\alpha_H,\delta_2$
and $\tau_2$ --- which is what Theorem~Theorem~4 of the thesis's remark at l.~1353--1355
requires of $\lambda_0$ and $\beta=\lambda_0/p$, since $r=\nu_1/\nu_0$ and
$\delta_2$ is built from $\delta$ and the data, while $c_2$ carries
$d(\Omega',\partial\Omega)$, $\|u\|_{L^\infty}$ and $\|u\|_{W^{1,p}}$ through
Theorem~3.  Only the norm $c_{20}$ is allowed to depend on $c_2$, and it does.
The four cited results are bound before $\alpha_2$; each mentions only $n$ and
$p$, both already bound, and each is universally quantified over its remaining
parameters, so none introduces a dependence of $\alpha_2$ or $\lambda_2$ on
anything outside the lists just given.
\end{helper}
\begin{proof}
Throughout, $B_t=B_t(x_0)$, $|E|$ is Lebesgue measure, $\omega_n=|B_1(0)|$, and
all non-integer powers are real powers of positive numbers.  Put
\[
\theta=\frac{\delta'}{n+p\alpha}>0 .
\]

\emph{Clause (A).}  Apply \noderef{CampanatoDecayFromBundle} with the parameters
$n,p,\alpha,\delta',\tau$.  Its parameter hypotheses are exactly the present
clause's: $n\ge2$, $p>1$, $0<\alpha\le1$, $\delta'>0$, $\tau>0$ and
$\tau<p\alpha\theta$.  It produces $\lambda_0\in(0,p)$ and $R_2>0$ with
$R_2\le1$, depending only on $n,p,\alpha,\delta'$ and $\tau$; feeding it the
present clause's $c_1,c_3>0$ then produces $c_{10}>0$ depending only on
$n,p,\alpha,\delta',\tau,c_1,c_3$.  These are the three constants clause (A)
announces, and the dependence lists are the ones its quantifier order
prescribes --- indeed the cited node binds $\lambda_0$ and $R_2$ before $c_1$
and $c_3$, which is more than clause (A) asks.

Now let $\varepsilon,\Omega,\Omega',u,g,R_1$ be as in clause (A).  The remaining
hypotheses of \noderef{CampanatoDecayFromBundle} hold verbatim:
$0\le\varepsilon\le1$; $\Omega$ open, $\Omega'$ open with
$\overline{\Omega'}\subseteq\Omega$, and $R_1>0$; $u\in W^{1,p}(\Omega)$ with
weak gradient $g$ (\noderef{MemW1pOn}); and the Morrey bound
(4.15) of the thesis.  For the comparison bundle, fix $x_0\in\Omega'$ and
$0<R\le R_1$ with $\overline B_R(x_0)\subseteq\Omega$: the present clause
supplies a comparison map $v\in W^{1,p}(B_R(x_0))$ with weak gradient $g_v$
(\noderef{MemW1pOn}) satisfying (4.16) of the thesis for $0<\rho\le R/2$ and
(4.17) of the thesis, and the cited node asks for exactly this, except that it
assumes (4.16) of the thesis only on the smaller range $0<\rho\le R/4$, which is
implied by the present clause's range because $R/4\le R/2$.  Its conclusion is
(4.19) of the thesis for every $x_0\in\Omega'$, every $R$ with
$0<R\le\min\{R_1,R_2\}$ and $\overline B_R(x_0)\subseteq\Omega$, and every
$\rho>0$ with $2\rho\le R^{1+\theta}$.  That is clause (A).


\emph{Step B1: the constants of clause (B).}  Let $r>0$ and $0<\alpha_H\le1$ be
given, together with the four cited results.  Apply
\noderef{GradientOscillationEstimate} (Theorem~2) with the parameters $n,p$ and
$r$ --- its hypotheses $2\le n$, $1<p$, $0<r$ hold, and its hypothesis
\noderef{CitedSobolevPoincareBall} and its two universal hypotheses, the cited
solvability result (C2) (\noderef{CitedDirichletW1p}) and the cited regularity
result (C1) (\noderef{CitedRegularityEpsPos}), are the present clause's.  Let
$C_2>0$ and $\alpha_2\in(0,1]$ be the constants it produces; they depend only on
$n,p$ and $r$.  This $\alpha_2$ is the exponent the clause announces.  Apply
\noderef{GradientSupEstimate} (Theorem~1) with the same $n,p,r$ and the same
three hypotheses, and let $C_1>0$ be its constant; and apply
\noderef{FrozenComparisonDecay} with $n,p,r$ and all four cited results, and let
$c_0>0$ be its constant.  All three depend only on $n,p$ and $r$.  Put
\[
\kappa_p=\max\bigl\{2^{\frac p2},\,2^{\,p-1}\bigr\}>0,
\qquad
c_3^{\sharp}=4^{\,p}\,\omega_n\,C_2^{\,p}\bigl(2^{\,n}C_1c_0\,\kappa_p\bigr)>0,
\]
both depending only on $n,p$ and $r$.  Each factor is used and checked where it
arises: $\kappa_p$ in the pointwise inequality of (B2b), the factor $2^{\,n}$ in
the sup bound of (B2b), and $4^{\,p}$ in the step of (B2c) that converts the
radius $R/4$ of Theorem~Theorem~2 of the thesis's application into the radius $R$ of
(4.16) of the thesis.  Now let $\delta_2>0$ and
$0<\tau_2<p\alpha_2\theta_2$ be given, $\theta_2=\delta_2/(n+p\alpha_2)$, and
let $c_2>0$ be given.

\emph{Step B2: the bundle, produced here.}  Fix data
$\varepsilon,\gamma_0,\gamma_1,\Omega,\Omega',\mathcal A,u,U,g,R_1$ satisfying
all the hypotheses of clause (B), and fix $x_0\in\Omega'$ and
$0<R\le\min\{R_1,1\}$.  (The restriction $R\le1$ costs nothing: Step B3 uses
the bundle only at radii $\le\min\{R_1,R_3\}$ with $R_3\le1$.)
Write $A^0=\mathcal A(x_0,U(x_0))$ and $\widetilde A=A^0-A^0(0)$.

\emph{(B2a) The comparison map.}  Apply the second part of
\noderef{FrozenComparisonDecay} --- the instance of Step B1 --- with the data
$\varepsilon$, $\alpha_H$, $\gamma_0$, $\gamma_1$, the field $A^0$, the centre
$x_0$, the radius $R$, the function $U$ and the field $g$.  Its hypotheses
hold: $0\le\varepsilon\le1$ and $0<\alpha_H\le1$ by hypothesis; $0<\gamma_0$,
$0<\gamma_1$ and $\gamma_1\le r\gamma_0$ by hypothesis; $0<R$; the field
$\widetilde A$ satisfies (Q1) of the thesis--(Q2) of the thesis on $\R^n$ with
$p,\varepsilon,\gamma_0,\gamma_1$ (\noderef{EllipticStructureQ12}), which is the
second hypothesis of clause (B) at $x_0$; $U\in W^{1,p}(B_R(x_0))$ with weak
gradient $g$ (\noderef{MemW1pOn}), by the hypothesis
$U\in W^{1,p}(\Omega)$ together with $B_R(x_0)\subseteq\overline
B_R(x_0)\subseteq\Omega$ and clause (1) of \noderef{WeakGradientRestrict}; and
$U\in C^{\alpha_H}(\overline B_R(x_0))$ (\noderef{MemHolderOn}), by the third
hypothesis of clause (B) at $x_0$ together with
$\overline B_R(x_0)\subseteq\overline B_{R_1}(x_0)$ and the fact that a function
of class $C^{\alpha_H}$ on a set is of class $C^{\alpha_H}$ on every subset,
with the same constant.  It produces $v,V,g_v$ with the frozen comparison bundle
$\mathrm{FrozenComparisonOn}(n,p,x_0,R,A^0,U,U,g,v,V,g_v)$
(\noderef{FrozenComparisonOn}) and the decay
\[
\int_{B_\rho(x_0)}\bigl(\varepsilon+|g_v|^2\bigr)^{\frac p2}dx
\le c_0\Bigl(\frac\rho R\Bigr)^{n}\int_{B_R(x_0)}\bigl(\varepsilon+|g_v|^2\bigr)^{\frac p2}dx
\qquad(0<\rho\le R).
\tag{4.11}
\]

\emph{(B2b) The sup bound, and $\Lambda$.}  Apply \noderef{GradientSupEstimate}
--- the instance of Step B1 --- with $\varepsilon,\gamma_0,\gamma_1$, the open
set $B_R(x_0)$, the field $\widetilde A$, the solution $v$ with weak gradient
$g_v$, the centre $x_0$ and the radius $R/2$.  Its hypotheses hold:
$0\le\varepsilon$; $0<\gamma_0$; $\gamma_1\le r\gamma_0$; $B_R(x_0)$ is open;
$0<R/2$; $\widetilde A(0)=0$ by construction; $\widetilde A$ satisfies
(Q1) of the thesis--(Q2) of the thesis on $B_R(x_0)$, a fortiori from $\R^n$;
$\overline B_{R/2}(x_0)\subseteq B_R(x_0)$; and $v$ is a weak solution of
$\operatorname{div}\widetilde A(\nabla v)=0$ in $B_R(x_0)$, which is the one
hypothesis that is not immediate and which we now prove, the first clause of the
bundle giving it only for $A^0$.

\emph{The operator shift.}  We claim first that
\[
\int_{B_R(x_0)}g_\varphi^{\,i}\,dx=0\qquad(i=1,\dots,n)
\tag{E3}
\]
for every $\varphi\in W^{1,p}(B_R(x_0))$ with weak gradient $g_\varphi$
(\noderef{MemW1pOn}) having compact support with
$\operatorname{supp}\varphi\subseteq B_R(x_0)$.  By
\noderef{WeakGradientZeroExtend} applied with $q=p$, the open set $\R^n$, the
ball $B_R(x_0)$ and the pair $(\varphi,g_\varphi)$ --- its hypotheses being
$1\le n$, $1\le p$, the openness of $\R^n$, $B_R(x_0)\subseteq\R^n$, the
membership just assumed and the compact support condition --- the pair
$(\varphi,\mathbf1_{B_R(x_0)}g_\varphi)$ satisfies $\varphi\in W^{1,p}(\R^n)$
with that weak gradient.  Let $\psi=\mathbf1_{\overline B_{R+2}(x_0)}*\rho_1$
with $\rho_1$ the standard mollifier of the Preamble at scale $1$: it is smooth
by the Preamble item on convolution with a smooth compactly supported kernel
(the indicator is locally integrable and $\rho_1$ is smooth with compact
support); it equals $1$ at every $x$ with $|x-x_0|\le R+1$, because then
$x-y\in\overline B_{R+2}(x_0)$ for every $y\in\operatorname{supp}\rho_1
\subseteq\overline B_1(0)$, so that $\psi(x)=\int\rho_1=1$; and it vanishes at
every $x$ with $|x-x_0|>R+3$, so $\operatorname{supp}\psi\subseteq
\overline B_{R+3}(x_0)$ is compact.  The defining identity of
\noderef{HasWeakGradientOn} on $\R^n$, which the pair
$(\varphi,\mathbf1_{B_R(x_0)}g_\varphi)$ satisfies by \noderef{MemW1pOn},
applied to the test function $\psi$ and the index $i$, reads
$\int_{\R^n}\varphi\,D_i\psi\,dx=-\int_{\R^n}\mathbf1_{B_R(x_0)}
g_\varphi^{\,i}\psi\,dx$.  Its left-hand side vanishes, because $\psi$ is
constant on the open ball $B_{R+1}(x_0)$, so $D_i\psi=0$ there, while
$\varphi=0$ off $\operatorname{supp}\varphi\subseteq B_R(x_0)\subseteq
B_{R+1}(x_0)$; and its right-hand side is $-\int_{B_R(x_0)}g_\varphi^{\,i}dx$,
because $\psi=1$ on $B_R(x_0)$.  This is (E3).

Now fix an admissible test pair $(\varphi,g_\varphi)$ for
\noderef{IsWeakSolutionOn} on $B_R(x_0)$, i.e.\ one with
$\varphi\in W^{1,p}(B_R(x_0))$ of weak gradient $g_\varphi$, compact support and
$\operatorname{supp}\varphi\subseteq B_R(x_0)$.  Write
\[
F=\sum_i\bigl(\widetilde A(g_v)\bigr)^ig_\varphi^{\,i},
\qquad
c=\sum_i\bigl(A^0(0)\bigr)^ig_\varphi^{\,i},
\]
so that, since $A^0(h)=\widetilde A(h)+A^0(0)$ for every $h$ and the field of
the present construction does not depend on the function value, the integrand
of the $A^0$-identity is $F+c$ and that of the $\widetilde A$-identity is $F$,
the right-hand sides of both identities being $0$.

\emph{$c$ is integrable with integral $0$.}  Each $g_\varphi^{\,i}$ is
integrable on $B_R(x_0)$: $|g_\varphi^{\,i}|\le|g_\varphi|\in L^p(B_R(x_0))$ by
\noderef{MemW1pOn} and $|B_R(x_0)|<\infty$, so H\"older's inequality on a set of
finite measure (Preamble) puts it in $L^1(B_R(x_0))$.  Hence $c$, a finite
linear combination of the $g_\varphi^{\,i}$ with the constant coefficients
$(A^0(0))^i$, is integrable on $B_R(x_0)$, and by linearity of the integral on
integrable functions together with (E3),
\[
\int_{B_R(x_0)}c\,dx=\sum_i\bigl(A^0(0)\bigr)^i\int_{B_R(x_0)}g_\varphi^{\,i}\,dx
=\sum_i\bigl(A^0(0)\bigr)^i\cdot0=0 .
\]

\emph{The three clauses of \noderef{IsWeakSolutionOn} match, and no
integrability of $F$ has to be assumed.}  That node asks for three things, and
each is compared separately.  (i) The membership clause
$v\in W^{1,p}_{\mathrm{loc}}(B_R(x_0))$ with weak gradient $g_v$
(\noderef{MemW1pLocOn}) does not mention the operator, so it is the same clause
for both fields.  (ii) The integrability clause holds for one field exactly when
it holds for the other: if $F$ is integrable on $B_R(x_0)$ then so is $F+c$,
being a sum of two integrable functions, and if $F$ is not integrable then
neither is $F+c$, since otherwise $F=(F+c)-c$ would be a difference of two
integrable functions.  (iii) The identity clause likewise holds for one exactly
when it holds for the other.  In the first case linearity of the integral and
the display above give
$\int_{B_R(x_0)}(F+c)\,dx=\int_{B_R(x_0)}F\,dx$, so the two integrals vanish
together.  In the second case neither $F$ nor $F+c$ is integrable, so both
integrals are $0$ by the convention that the integral of a non-integrable
function is $0$, and both identities hold --- though in that case the
integrability clause fails for both fields alike, so neither predicate is
asserted and (iii) is not needed.  As the admissible test pair was arbitrary,
all three clauses transfer, and $v$ is a weak solution for $\widetilde A$ in
$B_R(x_0)$.  Its conclusion is
\[
\bigl(\varepsilon+|g_v(x)|^2\bigr)^{\frac p2}
\le\frac{2^{\,n}C_1}{R^{\,n}}\int_{B_{R/2}(x_0)}\bigl(\varepsilon+|g_v|^2\bigr)^{\frac p2}dy
\qquad\text{for a.e. }x\in B_{R/4}(x_0).
\]
By (4.11) of the thesis at $\rho=R/2$, which gives
$\int_{B_{R/2}(x_0)}(\varepsilon+|g_v|^2)^{p/2}\le c_02^{-n}
\int_{B_R(x_0)}(\varepsilon+|g_v|^2)^{p/2}$, the right-hand side is at most
$2^{\,n}C_1\cdot c_02^{-n}R^{-n}\int_{B_R(x_0)}(\varepsilon+|g_v|^2)^{p/2}dx
=C_1c_0R^{-n}\int_{B_R(x_0)}(\varepsilon+|g_v|^2)^{p/2}dx$; the factor $2^{-n}$
is the one (4.11) of the thesis actually supplies at $\rho=R/2$ and it is kept.

Next, the pointwise comparison of $g_v$ with $\nabla u$.  From
$|g_v|\le|\nabla u|+|\nabla u-g_v|$ and $(X+Y)^2\le2(X^2+Y^2)$,
\[
1+|g_v|^2\le1+2|\nabla u|^2+2|\nabla u-g_v|^2
\le2\bigl[(1+|\nabla u|^2)+|\nabla u-g_v|^2\bigr],
\]
and raising to the power $p/2$ and using
$(X+Y)^{p/2}\le\max\{1,2^{\frac p2-1}\}(X^{p/2}+Y^{p/2})$ for $X,Y\ge0$ ---
subadditivity of $t\mapsto t^{p/2}$ when $p\le2$, convexity when $p\ge2$, which
is (E1) of Step 0 with the exponent $p/2$ in place of $p$ --- gives, with
$X=1+|\nabla u|^2$ and $Y=|\nabla u-g_v|^2$ so that $Y^{p/2}=|\nabla u-g_v|^p$,
\[
(\varepsilon+|g_v|^2)^{\frac p2}\le(1+|g_v|^2)^{\frac p2}
\le2^{\frac p2}\max\{1,2^{\frac p2-1}\}
\bigl[(1+|\nabla u|^2)^{\frac p2}+|\nabla u-g_v|^{p}\bigr]
=\kappa_p\bigl[(1+|\nabla u|^2)^{\frac p2}+|\nabla u-g_v|^{p}\bigr],
\]
using $(\varepsilon+t^2)^{p/2}\le(1+t^2)^{p/2}$ for $\varepsilon\le1$ and
$2^{p/2}\max\{1,2^{p/2-1}\}=\max\{2^{p/2},2^{\,p-1}\}=\kappa_p$.  Integrating
over $B_R(x_0)$,
$\int_{B_R}(\varepsilon+|g_v|^2)^{p/2}\le\kappa_p\{\int_{B_R}(1+|\nabla
u|^2)^{p/2}+\int_{B_R}|\nabla u-g_v|^{p}\}$, and the second term is at most
$R^{\delta_2+n}\Lambda(x_0,R)=R^{\delta_2}\int_{B_R}(1+|\nabla u|^2)^{p/2}
\le\int_{B_R}(1+|\nabla u|^2)^{p/2}$ by the fifth hypothesis of clause (B),
applied to the bundle of (B2a), and $R^{\delta_2}\le1$ because $R\le1$ and
$\delta_2>0$.  Hence
$\int_{B_R}(\varepsilon+|g_v|^2)^{p/2}\le2\kappa_pR^{\,n}\Lambda(x_0,R)$, and
the sup bound becomes, for a.e.\ $x\in B_{R/4}(x_0)$,
\[
\bigl(\varepsilon+|g_v(x)|^2\bigr)^{\frac p2}
\le2C_1c_0\kappa_p\,\Lambda(x_0,R)
\le S^{p},
\qquad
S:=\Bigl(2^{\,n}C_1c_0\kappa_p\,\Lambda(x_0,R)\Bigr)^{1/p},
\]
the last inequality because $2\le2^{\,n}$ for $n\ge2$ (here $n\ge2$ by
hypothesis, and $n\ge1$ would do).
In particular $|g_v(x)|\le S$ almost everywhere on $B_{R/4}(x_0)$: raising the
display to the power $1/p>0$, which is monotone on $[0,\infty)$, gives
$(\varepsilon+|g_v(x)|^2)^{1/2}\le S$, and $|g_v(x)|\le(\varepsilon+|g_v(x)|^2)^{1/2}$
because $\varepsilon\ge0$.

\emph{(B2c) The oscillation estimate (4.16) of the thesis.}  Apply
\noderef{GradientOscillationEstimate} --- the instance of Step B1 --- with
$\varepsilon,\gamma_0,\gamma_1$, the bound $S$ of (B2b), the open set
$B_R(x_0)$ --- the \emph{same} open set as in (B2b), so that no restriction of
the weak-solution property to a smaller ball is needed --- the field
$\widetilde A$, the solution $v$ with gradient $g_v$, the centre $x_0$ and the
radii $R/4$ and $\rho\le R/4$.  Its hypotheses are exactly those of (B2b) ---
$0\le\varepsilon$, $0<\gamma_0$, $\gamma_1\le r\gamma_0$, the openness of
$B_R(x_0)$, (Q1) of the thesis--(Q2) of the thesis for $\widetilde A$ on $B_R(x_0)$ and
the weak-solution property of $v$ for $\widetilde A$ on $B_R(x_0)$, proved
there --- together with $0<\rho\le R/4$,
$\overline B_{R/4}(x_0)\subseteq B_R(x_0)$ and $|g_v|\le S$ almost everywhere on
$B_{R/4}(x_0)$, which is the last display of (B2b).  Its conclusion is
$\|g_v(x)-g_v(y)\|\le C_2S(\rho/(R/4))^{\alpha_2}=C_2S(4\rho/R)^{\alpha_2}$ for
almost every pair $x,y\in B_\rho(x_0)$.  The field $g_v$ is integrable on
$B_\rho(x_0)$ --- $|g_v|\le S$ a.e.\ there, since
$B_\rho(x_0)\subseteq B_{R/4}(x_0)$, and $|B_\rho(x_0)|<\infty$ --- so its mean
$(g_v)_{B_\rho(x_0)}$ is defined, and for almost every $x\in B_\rho(x_0)$
\[
\bigl\|g_v(x)-(g_v)_{B_\rho(x_0)}\bigr\|
=\Bigl\|\frac1{|B_\rho|}\int_{B_\rho}\bigl(g_v(x)-g_v(y)\bigr)dy\Bigr\|
\le\frac1{|B_\rho|}\int_{B_\rho}\bigl\|g_v(x)-g_v(y)\bigr\|\,dy
\le C_2S\Bigl(\frac{4\rho}R\Bigr)^{\alpha_2},
\]
by the norm inequality for the Bochner integral over a general measure
(Mathlib's \texttt{MeasureTheory.norm\_integral\_le\_integral\_norm}) and the
a.e.\ bound just obtained, valid for a.e.\ $y$ at each such $x$.  Raising to the
power $p$ and integrating over $B_\rho(x_0)$, whose measure is
$\omega_n\rho^{\,n}$ (Preamble),
\[
\int_{B_\rho(x_0)}\bigl|g_v-(g_v)_{B_\rho(x_0)}\bigr|^p dx
\le\omega_nC_2^{\,p}4^{\,p\alpha_2}S^{p}\Bigl(\frac\rho R\Bigr)^{p\alpha_2}\rho^{\,n}
\le c_3^{\sharp}\Bigl(\frac\rho R\Bigr)^{p\alpha_2}\rho^{\,n}\Lambda(x_0,R)
\]
for $0<\rho\le R/4$, with $c_3^{\sharp}$ of Step B1: the first inequality is
$S^{p}=2^{\,n}C_1c_0\kappa_p\Lambda(x_0,R)$ from (B2b), and the second is
$4^{\,p\alpha_2}\le4^{\,p}$, which holds because $0<\alpha_2\le1$ and $4^p\ge1$,
together with the definition
$c_3^{\sharp}=4^{\,p}\omega_nC_2^{\,p}(2^{\,n}C_1c_0\kappa_p)$ --- the factor
$4^{\,p}$ is \emph{inside} $c_3^{\sharp}$, not absorbed by an inequality
$4^{p\alpha_2}\le1$, which is false.

Together with the fifth hypothesis of clause (B), this is exactly the pair
(4.16) of the thesis--(4.17) of the thesis that \noderef{CampanatoDecayFromBundle}
assumes, for the field $g_v$ just produced, at the exponent $\alpha_2$, the
constant $c_3^{\sharp}$ and the Morrey data $\tau_2,c_2$.  The ranges match
exactly and no radius conversion is needed: that node assumes (4.16) of the thesis
on $0<\rho\le R/4$, which is precisely the range produced here, and the
restriction to $R/4$ is the reason it is stated there on that range rather than
on the thesis' $\rho\le R/2$.  The $W^{1,p}$ membership that node also assumes
is the second clause of the bundle
$\mathrm{FrozenComparisonOn}(n,p,x_0,R,A^0,U,U,g,v,V,g_v)$ obtained in (B2a),
namely $v\in W^{1,p}(B_R(x_0))$ with weak gradient $g_v$
(\noderef{FrozenComparisonOn}, \noderef{MemW1pOn}).

\emph{Step B3: the conclusion of clause (B).}  Apply
\noderef{CampanatoDecayFromBundle} with the parameters
$n,p,\alpha_2,\delta_2,\tau_2$ in place of $n,p,\alpha,\delta',\tau$.  Its
parameter hypotheses hold: $n\ge2$ and $p>1$ are the present clause's;
$0<\alpha_2\le1$ by Step B1; and $\delta_2>0$, $\tau_2>0$ and
$\tau_2<p\alpha_2\theta_2$ with $\theta_2=\delta_2/(n+p\alpha_2)$ are
hypotheses of clause (B).  It produces $\lambda_2\in(0,p)$ and $R_3>0$ with
$R_3\le1$, depending only on $n,p,\alpha_2,\delta_2$ and $\tau_2$.  These are
the $\lambda_2$ and $R_3$ the clause announces; since $\alpha_2$ is a function
of $n,p$ and $r$ alone (Step B1), they depend only on $n,p,r,\delta_2$ and
$\tau_2$, which is what Theorem~Theorem~4 of the thesis's remark at l.~1353--1355 requires.
They are produced before $c_2$ is given, which is the clause's quantifier
order; this is available because the cited node binds its $\lambda_0$ and $R_2$
before its $c_1$ and $c_3$.

Now let $c_2>0$ be given.  Feeding the cited node the constants $c_2$ and
$c_3^{\sharp}$ --- both positive, $c_3^{\sharp}$ by Step B1 --- produces
$c_{20}>0$ depending only on $n,p,\alpha_2,\delta_2,\tau_2,c_2$ and
$c_3^{\sharp}$, hence on $n,p,r,\delta_2,\tau_2$ and $c_2$.  This is the
$c_{20}$ the clause announces, and it is the only constant allowed to depend on
$c_2$.

Finally fix the data $\varepsilon,\gamma_0,\gamma_1,\Omega,\Omega',\mathcal
A,u,U,g,R_1$ of clause (B), and apply the cited node at the radius parameter
$R_1^{\flat}=\min\{R_1,1\}>0$ in place of $R_1$.  (Its $R_1$ is universally
quantified, so this is an instantiation and not a restriction; $R_1^{\flat}$ is
used because Step B2 produces the bundle only for $R\le\min\{R_1,1\}$, and
nothing is lost because the conclusion will be needed only for
$R\le R_3\le1$.)  The remaining hypotheses hold: $0\le\varepsilon\le1$;
$\Omega$ open, $\Omega'$ open with $\overline{\Omega'}\subseteq\Omega$; and
$u\in W^{1,p}(\Omega)$ with weak gradient $g$ (\noderef{MemW1pOn}), which is a
hypothesis of clause (B).  The Morrey bound $\Lambda(x_0,R)\le c_2R^{-\tau_2}$
for $x_0\in\Omega'$ and $0<R\le R_1^{\flat}$ is the fourth hypothesis of clause
(B), whose range $0<R\le R_1$ contains this one; it is (4.15) of the thesis at the
substituted constants.  The comparison bundle is what Step B2 produces: for
every $x_0\in\Omega'$ and every $0<R\le R_1^{\flat}=\min\{R_1,1\}$ with
$\overline B_R(x_0)\subseteq\Omega$ --- which is exactly the range Step B2 was
run on --- Step B2 yields $v,V,g_v$ with $v\in W^{1,p}(B_R(x_0))$ and weak
gradient $g_v$ (the second clause of its bundle), the oscillation estimate
(4.16) of the thesis at the exponent $\alpha_2$ and the constant $c_3^{\sharp}$ for
$0<\rho\le R/4$ by (B2c), and the smallness (4.17) of the thesis at $\delta_2$,
which is the fifth hypothesis of clause (B) applied to that bundle.

Its conclusion is: for every $x_0\in\Omega'$, every $R$ with
$0<R\le\min\{R_1^{\flat},R_3\}$ and $\overline B_R(x_0)\subseteq\Omega$, and
every $\rho>0$ with $2\rho\le R^{1+\theta_2}$,
\[
\int_{B_\rho(x_0)}\bigl|\nabla u-(\nabla u)_{B_\rho(x_0)}\bigr|^p dx
\le c_{20}\,\rho^{\,n+\lambda_2}.
\]
Here $\min\{R_1^{\flat},R_3\}=\min\{R_1,1,R_3\}=\min\{R_1,R_3\}$ because
$R_3\le1$, so the range of $R$ is the one clause (B) claims.
The containment $\overline B_R(x_0)\subseteq\Omega$ is automatic here from the
first hypothesis of clause (B), which asserts it for every $x_0\in\Omega'$ and
every $0<R\le R_1$, and $R\le\min\{R_1,R_3\}\le R_1$.  So the displayed estimate
holds for every $x_0\in\Omega'$, every $0<R\le\min\{R_1,R_3\}$ and every
$\rho>0$ with $2\rho\le R^{1+\theta_2}$, which is clause (B).  The exponent
$\alpha_H$ enters only through the production of the bundle in Step B2 and not
through any constant, so listing it among the dependences of $\alpha_2$,
$\lambda_2$ and $R_3$ over-reports rather than under-reports them.
\end{proof}
